% Incorporación ampliativa DJ-05 del inventario REV05.
% Residencia: artículo III, después de inversión cúbica y recuperación angular.
% Propietario: output/REPARACION_ENTREGA_MULTIPLATAFORMA_20260916/ENTREGA/ES/
% PAQUETES/03_VACIO_ELECTROMAGNETICO_ES/payload/spanish_source/manuscrito/
% sections/04_respuesta_constitutiva.tex, líneas 14--65 y 139--186.
% Estatuto: FORMALIZACION_NUEVA de la diferenciación y del dominio imagen;
% la respuesta, su monotonía y su inversa cúbica son RESULTADO_RECUPERADO.

\section[Inversion différentielle de la réponse constitutive]{Inversion différentielle et domaine de la réponse constitutive à deux feuillets}
\label{jc:sec:constitutiva}

Les canaux angulaires engendrés par le calendrier TPK déterminent le transport positif
\[
T=e^{-(x+y)}P_++e^{-(x-y)}P_-,\qquad
\Omega=\{(x,y)\in\mathbb R^2:x>|y|\}.
\]
Ici, $P_+$ et $P_-$ sont les projecteurs étiquetés des deux feuillets. Les incidences $90$ et $120$ agissent sur ce même transport et produisent la réponse
\[
\mathcal V=(I-T^{90})(I-T^{120})^{-1}=r_+P_++r_-P_-.
\]
La factorisation avec $S=T^{30}$ détermine
\begin{equation}
f(s)=\frac{1+s+s^2}{1+s+s^2+s^3},\qquad
r_\pm=f\bigl(e^{-30(x\pm y)}\bigr).
\label{jc:eq:respuesta}
\end{equation}
Le facteur $30$ conserve l’unité commune des incidences $90=3\cdot30$ et $120=4\cdot30$. Cette section reçoit les coordonnées angulaires précédemment produites ; elle étudie la régularité et l’inversion de leur lecture constitutive.

\subsection{Sensibilité spectrale et coordonnées logarithmiques}

Définissons, pour $t>0$,
\[
g(t)=\log f(e^{-30t}),\qquad \sigma(t)=g'(t),
\]
et les coordonnées de quotient et de produit inverse
\begin{align}
B(x,y)&=\log\frac{r_-}{r_+}=g(x-y)-g(x+y),\\
V(x,y)&=\log\frac1{r_-r_+}=-g(x-y)-g(x+y).
\label{jc:eq:bv}
\end{align}
Les lettres $B,V$ désignent ici deux coordonnées scalaires de la réponse ; les projecteurs de feuillet demeurent dans l’opérateur $\mathcal V$.

\begin{lemma}[Monotonie régulière du lecteur spectral]
L’application $g$ est une bijection analytique strictement croissante de $(0,\infty)$ dans $(-\log(4/3),0)$. Avec $s=e^{-30t}$,
\begin{equation}
\sigma(t)=
\frac{30s^3(3+2s+s^2)}{(1+s+s^2)(1+s+s^2+s^3)}>0.
\label{jc:eq:sigma}
\end{equation}
\end{lemma}
\begin{proof}
La dérivation du quotient donne
\[
f'(s)=-\frac{s^2(3+2s+s^2)}{(1+s+s^2+s^3)^2}<0,
\qquad 0<s<1.
\]
La règle de dérivation composée produit $g'(t)=-30s f'(s)/f(s)$, qui est \eqref{jc:eq:sigma}. Les dénominateurs sont positifs sur le domaine. Les limites $f(1)=3/4$ et $f(0)=1$ donnent $g(0^+)=-\log(4/3)$ et $g(+\infty)=0$. La monotonie, la continuité et ces limites démontrent la bijectivité. La positivité de la dérivée garantit la régularité analytique de l’inverse sur l’intervalle ouvert.
\end{proof}

\begin{theorem}[Jacobien constitutif et récupération locale]
\label{jc:thm:jacobiano} Soient $a=\sigma(x-y)$ et $b=\sigma(x+y)$. L’application $\Phi(x,y)=(B(x,y),V(x,y))$ satisfait
\begin{equation}
D\Phi=
\begin{pmatrix}a-b&-a-b\\-a-b&a-b\end{pmatrix},\qquad
\det D\Phi=-4ab<0.
\label{jc:eq:jacobiano}
\end{equation}
En particulier, elle est localement inversible sur tout $\Omega$ et les identités différentielles $B_x=V_y$ et $B_y=V_x$ sont satisfaites.
\end{theorem}
\begin{proof}
La dérivation directe de \eqref{jc:eq:bv} donne $B_x=a-b$, $B_y=-a-b$, $V_x=-a-b$ et $V_y=a-b$. Le déterminant est $(a-b)^2-(a+b)^2=-4ab$, strictement négatif par \eqref{jc:eq:sigma}. Le théorème d’inversion locale s’applique en chaque point du domaine. Les deux identités différentielles résultent des entrées égales de la matrice.
\end{proof}

\subsection{Domaine image et reconstruction globale}

\begin{theorem}[Coordonnées globales de la réponse normalisée]
\label{jc:thm:global} Soit $L=\log(4/3)$. L’application $\Phi$ est un difféomorphisme analytique de $\Omega$ sur
\begin{equation}
\mathcal D=\{(B,V):0<V-B<2L,\quad0<V+B<2L\}.
\label{jc:eq:diamante}
\end{equation}
Si $\psi=g^{-1}$, la reconstruction est donnée par
\begin{align}
x&=\frac12\left[\psi\!\left(\frac{B-V}{2}\right)
                   +\psi\!\left(-\frac{B+V}{2}\right)\right],\\
y&=\frac12\left[\psi\!\left(-\frac{B+V}{2}\right)
                   -\psi\!\left(\frac{B-V}{2}\right)\right].
\label{jc:eq:inversa}
\end{align}
La chambre $x>y>0$ correspond à la moitié $B<0$ de $\mathcal D$.
\end{theorem}
\begin{proof}
Le changement $(x,y)\mapsto(t_-,t_+)=(x-y,x+y)$ est une bijection linéaire entre $\Omega$ et $(0,\infty)^2$. Appliquer $g$ aux deux coordonnées donne une bijection analytique sur $(-L,0)^2$. Si $u=g(t_-)$ et $v=g(t_+)$, alors $B=u-v$ et $V=-u-v$. Son inverse linéaire est $u=(B-V)/2$, $v=-(B+V)/2$. Les conditions $-L<u,v<0$ équivalent exactement à \eqref{jc:eq:diamante}. La composition de ces trois bijections donne \eqref{jc:eq:inversa} et démontre l’assertion globale. Enfin, $y>0$ équivaut à $t_+>t_-$ ; la monotonie de $g$ la transforme en $v>u$, soit $B<0$.
\end{proof}

La fonction $\psi$ est évaluée par l’inverse cubique du lecteur déjà construit. Pour $u\in(-L,0)$, on prend $r=e^u$ et l’unique racine $s\in(0,1)$ de
\[
rs^3+(r-1)(1+s+s^2)=0.
\]
Alors $\psi(u)=-(\log s)/30$. La reconstruction différentielle et la reconstruction cubique sont donc deux expressions de la même inversion constitutive.

\begin{corollary}[Conjugaison des feuillets]
L’involution $(x,y)\mapsto(x,-y)$ agit sur la réponse par $(B,V)\mapsto(-B,V)$. Elle conserve le produit $r_+r_-$ et inverse le quotient $r_-/r_+$.
\end{corollary}
\begin{proof}
L’involution échange $x-y$ avec $x+y$ et, par conséquent, $r_-$ avec $r_+$. La formule \eqref{jc:eq:bv} fournit le résultat.
\end{proof}

\subsection{Stabilité locale de la récupération}

\begin{proposition}[Valeurs singulières et borne de sensibilité]
Les valeurs singulières de $D\Phi$ sont $2a$ et $2b$. En particulier,
\[
\bigl\|(D\Phi)^{-1}\bigr\|_2
=\frac1{2\min\{\sigma(x-y),\sigma(x+y)\}}.
\]
Pour $0<\delta<M<\infty$, la récupération différentielle est uniformément bornée sur $\Omega_{\delta,M}=\{(x,y):\delta\le x-y,x+y\le M\}$.
\end{proposition}
\begin{proof}
La matrice symétrique \eqref{jc:eq:jacobiano} possède les vecteurs orthogonaux $(1,-1)$ et $(1,1)$ comme vecteurs propres, de valeurs propres $2a$ et $-2b$. Les valeurs singulières sont leurs valeurs absolues. La norme spectrale de l’inverse est l’inverse de la plus petite valeur singulière. La fonction continue positive $\sigma$ atteint un minimum positif sur $[\delta,M]$, ce qui fournit la borne uniforme.
\end{proof}

\begin{example}[Récupération exacte de deux canaux rationnels]
Prenons $e^{-30(x-y)}=1/2$ et $e^{-30(x+y)}=1/4$. Alors
\[
x=\frac{\log2}{20},\qquad y=\frac{\log2}{60},\qquad
r_-=\frac{14}{15},\qquad r_+=\frac{84}{85}.
\]
Les coordonnées constitutives sont
\[
B=\log\frac{17}{18},\qquad V=\log\frac{425}{392}.
\]
L’inversion cubique récupère $1/2$ et $1/4$, et les formules \eqref{jc:eq:inversa} restituent $x,y$. En ce point,
\[
a=\frac{34}{7},\qquad b=\frac{114}{119},\qquad
\det D\Phi=-\frac{15504}{833}.
\]
Tous les coefficients rationnels de l’exemple s’obtiennent en évaluant la fonction constitutive et sa dérivée ; l’exemple vérifie l’inversion sans modifier le générateur des canaux.
\end{example}

La transformation $(x,y)\leftrightarrow(B,V)$ décrit exactement l’information angulaire accessible par les deux lectures normalisées. Les réalisations dimensionnelles de permittivité, de perméabilité, d’impédance et de vitesse sont ensuite appliquées avec les droites dimensionnelles et leurs sections déclarées. La signature différentielle précédente est une propriété de cette application de coordonnées constitutives et conserve ce domaine d’interprétation.
